\documentclass[11pt]{amsart}
\usepackage[margin=1.1in]{geometry}
\usepackage{amsmath,amssymb,amsthm,mathtools}
\usepackage[T1]{fontenc}
\usepackage{lmodern}
\usepackage{microtype}
\usepackage{enumitem}
\usepackage{booktabs}
\usepackage[dvipsnames]{xcolor}
\usepackage[colorlinks=true,linkcolor=blue!50!black,citecolor=blue!50!black,urlcolor=blue!50!black]{hyperref}
\usepackage[capitalise,nameinlink]{cleveref}
\crefname{step}{Step}{Steps}
\crefname{status}{Status}{Statuses}

\newtheorem{theorem}{Theorem}[section]
\newtheorem{proposition}[theorem]{Proposition}
\newtheorem{lemma}[theorem]{Lemma}
\newtheorem{corollary}[theorem]{Corollary}
\theoremstyle{definition}
\newtheorem{definition}[theorem]{Definition}
\newtheorem{step}[theorem]{Step}
\newtheorem{status}[theorem]{Status}
\theoremstyle{remark}
\newtheorem{remark}[theorem]{Remark}

\newcommand{\QQ}{\mathbb{Q}}
\newcommand{\ZZ}{\mathbb{Z}}
\newcommand{\CC}{\mathbb{C}}
\newcommand{\cO}{\mathcal{O}}
\newcommand{\cI}{\mathcal{I}}
\newcommand{\cF}{\mathcal{F}}
\newcommand{\ch}{\operatorname{ch}}
\newcommand{\Ext}{\operatorname{Ext}}
\newcommand{\cExt}{\mathcal{E}xt}
\newcommand{\Hdg}{\operatorname{Hdg}}
\newcommand{\NS}{\operatorname{NS}}
\newcommand{\Nm}{\operatorname{Nm}}
\newcommand{\id}{\operatorname{id}}
\newcommand{\At}{\mathrm{At}}

\title{Hilbert--Burch Resolutions of Secant Supports on Abelian Fourfolds}
\author{Deep Bhattacharjee}
\author{Priyabrata Mandal}
\author{Ushashi Bhattacharya}
\date{Working note, 29 September 2026; corrected version of the route note of
the same day}

\begin{document}

\begin{abstract}
We prove that the ideal sheaf of the support of a secant object on a
principally polarised abelian fourfold has no Hilbert--Burch resolution of
rank one, and none of rank two at the discriminants $1$ and $5$ when its
Chern classes are polynomials in the polarisation, and we correct the route
note that proposed these supports.
\end{abstract}

\maketitle

\section{What this note does}\label{sec:status}

\begin{status}\label{st:proves}
This note proves no case of the Hodge conjecture. It writes the secant route
of \cite[Section~18]{BMB} on one principally polarised abelian fourfold as a
list of steps, records which constructions are excluded, and proves one new
exclusion, \Cref{prop:lowrank}: the Hilbert--Burch resolution of a secant
support has rank $r\ge2$, and $r\ge3$ at $d=1,5$ when its Chern classes are
polynomials in $\Theta$. That proposition is also in the paper
\cite[Proposition~18.60]{BMB}. The numerical statements of
\Cref{sec:program,sec:lemma} are checked in exact arithmetic
(\Cref{sec:checks}); every other statement quoted here is proved in
\cite{BMB} or in the papers it cites, at the place given.
\end{status}

\subsection*{The target}
Markman proved that the Weil classes are algebraic on the abelian sixfolds of
split Weil type for every imaginary quadratic field \cite{Mar25a,Mar25b}, and
deduced it for all abelian fourfolds of Weil type, for every field and every
discriminant, through a product argument of Schoen
\cite[Proposition~10]{Sch98}; the Hodge conjecture for abelian fourfolds
follows \cite[Section~1]{BMB}. The fourfolds of Weil type are therefore not a
target. The secant route starts from a principally polarised abelian fourfold
$X$ and a secant object $F$ on it, and aims at the Weil classes
$\mathbf{W}(K,4,\delta)$ on the abelian eightfolds $X\times\widehat X$ of
Weil type, $K=\QQ(\sqrt{-d})$ \cite[Theorem~18.29(iv),
Remark~18.69]{BMB}. Besides an object on $X$ it needs an affirmative answer
to \cite[Question~11.4]{Mar25b}.

\subsection*{Where the route sits}
The closure theorem \cite[Theorem~20.58]{BMB} lists five minimal sufficient
sets of open statements,
\[
  \{\mathrm{F3}\},\quad\{\mathrm{L},\mathrm{M}\},\quad
  \{\mathrm{L},\mathrm{F3}'\},\quad\{\mathrm{V},\mathrm{F3}'\},\quad
  \{\mathrm{P2}_{\mathrm{split}},\mathrm{F2},\mathrm{F3}'\}.
\]
Statement (F1) of the paper is the propagation statement, and it is needed
only in the form $\mathrm{P2}_{\mathrm{split}}$, for the split families of
the CM fields of degree at least four: (P2) for an imaginary quadratic field
and for a nontrivial discriminant is derived by descent
\cite[Proposition~20.13]{BMB}. The secant
route lies in no minimal sufficient set \cite[Theorem~20.58(iv)]{BMB}:
granting both of its open demands gives $\mathbf{W}(K,n,\delta)$ for
imaginary quadratic $K$ and leaves the CM fields of degree at least four
outside. Statement (F3), the conjecture for the varieties that are not
abelian, is equivalent to the Hodge conjecture
\cite[Proposition~20.35]{BMB}; (F3$'$), the conjecture modulo abelian
varieties \cite[Definition~20.36]{BMB}, is weaker; and (F2) is equivalent to
the Hodge conjecture for abelian varieties \cite[Proposition~21.49]{BMB}.
Every element of a minimal sufficient set other than
$\mathrm{P2}_{\mathrm{split}}$ is a consequence of the Hodge conjecture, and
$\{\mathrm{F3}\}$ and $\{\mathrm{L},\mathrm{M}\}$ are each equivalent to it
\cite[Theorem~20.58(vi)]{BMB}.

\section{The reduction to one object}\label{sec:reduction}

Let $(X,\Theta)$ be a principally polarised abelian fourfold, $d>0$, and $F$
a perfect complex on $X$ with $\ch(F)=a\,u_{\Theta}+b\,v_{\Theta}$, where
$u_{\Theta},v_{\Theta}$ span the secant plane \cite[Section~18]{BMB}. No
action of $K$ on $X$ is assumed: the field $K=\QQ(\sqrt{-d})$ enters through
$\ch(F)$, and the Weil structure lives on $X\times\widehat X$
\cite[Proposition~18.25(v), Theorem~18.29(iv)]{BMB}.

\begin{theorem}[{\cite[Theorem~18.29(ii), (iii)]{BMB}}]\label{thm:factor}
The classes of $HT^{2}(X)$ preserving $\ch(F)$ to first order are the ten
polarised first order deformations of $(X,\Theta)$ and the six Poisson
directions compensated in $H^{0,2}$. The object $F$ satisfies the weakened
criterion if and only if
\begin{enumerate}[label=\textup{(\alph*)},nosep]
\item $v\lrcorner\At(F)=0$ in $\Ext^{2}(F,F)$ for every $v\in H^{1}(T_{X})$
  with $v\lrcorner\Theta=0$, that is, $F$ extends to first order along every
  polarised deformation of $(X,\Theta)$; and
\item $\At(F)^{2}=-d\,\Theta^{2}\cdot\id_{F}$ in
  $\Ext^{2}(F,F\otimes\Omega^{2}_{X})$.
\end{enumerate}
The trace of \textup{(b)} is automatic.
\end{theorem}

For Markman's object, built from translates of $\Theta$, condition (a) holds
and (b) is the whole condition \cite[Corollary~18.30]{BMB}. For any other
candidate both are needed, and (a) runs over all ten polarised directions.

\section{The local dichotomy and the numerical targets}\label{sec:targets}

Let $F=\cI_{Z}(b\Theta)$ with $Z$ of codimension two. If the local ring
$\cO_{Z,p}$ is Cohen--Macaulay, then $\cExt^{q}(\cI_{Z},\cI_{Z})_{p}=0$ for
$q\ge2$ \cite[Theorem~18.48]{BMB}. If $Z$ fails to be Cohen--Macaulay at a
point of one of the types of \cite[Remark~18.49]{BMB}, then $F$ fails the
criterion; if $Z$ is Cohen--Macaulay, the criterion is the vanishing of the
image of $A_{\theta}A_{\theta'}$ in $H^{1}(Z,\cExt^{1}(\cI_{Z},\cI_{Z}))$ for
all $\theta,\theta'\in H^{0}(T_{X})$, and when $Z$ is smooth
$\cExt^{1}(\cI_{Z},\cI_{Z})=N_{Z/X}$ \cite[Corollary~18.50]{BMB}.

For a smooth support $S$, Grothendieck--Riemann--Roch with trivial Todd class
forces \cite[Theorem~18.40]{BMB}
\[
  [S]=N\Theta^{2},\qquad N=\tfrac12(b^{2}+d),\qquad
  \chi(\cO_{S})=(b^{2}+d)(3b^{2}-d),
\]
\[
  K_{S}^{2}=3(b^{2}+d)(7b^{2}-d),\qquad e(S)=3(b^{2}+d)(5b^{2}-3d),\qquad
  K_{S}\cdot\lambda=32bN,\qquad\lambda^{2}=24N,
\]
with $\lambda=\Theta|_{S}$, and the Bogomolov--Miyaoka--Yau inequality is
exactly $d\le b^{2}$. At $b=3$ with $d$ squarefree and $N$ integral this
leaves $d\in\{1,3,5,7\}$ \cite[Corollary~18.41]{BMB}, with the invariants of
\Cref{tab:targets}.

\begin{table}[ht]
\centering
\begin{tabular}{@{}ccrrrrr@{}}
\toprule
$d$ & $N$ & $\chi(\cO_S)$ & $K_S^2$ & $e(S)$ & $K_S\cdot\lambda$ & $\lambda^2$\\
\midrule
1 & 5 & 260 & 1860 & 1260 & 480 & 120\\
3 & 6 & 288 & 2160 & 1296 & 576 & 144\\
5 & 7 & 308 & 2436 & 1260 & 672 & 168\\
7 & 8 & 320 & 2688 & 1152 & 768 & 192\\
\bottomrule
\end{tabular}
\caption{The forced invariants of a smooth support at $b=3$.}
\label{tab:targets}
\end{table}

\section{What is excluded}\label{sec:excluded}

A candidate in any of the following classes fails before condition (b) is
reached.
\begin{enumerate}[label=(\roman*),leftmargin=2.2em]
\item Complete intersections of two divisors \cite[Theorem~18.54]{BMB}.
\item Zero loci of sections of vector bundles of rank two, and more generally
  every Hilbert--Burch resolution of rank one (\Cref{prop:lowrank}(i)).
\item Cones of maps between sums of powers of an ample line bundle with class
  in $\QQ\Theta$, whatever the rank \cite[Theorem~18.63,
  Corollary~18.68(i)]{BMB}.
\item Cones of maps between sums of simple isogeny blocks, that is, of simple
  semi-homogeneous bundles \cite[Theorem~18.66, Remark~18.67]{BMB}; these
  include all line bundles.
\item For a smooth support, resolutions whose two terms both have classes in
  the lattice generated by line bundles \cite[Corollary~18.58]{BMB}.
\item Markman's explicit complex from $N$ translates of $\Theta$
  \cite[Theorem~18.34]{BMB}.
\item Hilbert--Burch resolutions of rank two at $d\equiv1\pmod4$, when the
  Chern classes are polynomials in $\Theta$ (\Cref{prop:lowrank}(ii)).
\end{enumerate}

\begin{proposition}[{\cite[Proposition~18.60]{BMB}}]\label{prop:lowrank}
Let $(X,\Theta)$ be a principally polarised abelian fourfold, $b\ge1$, $d>0$,
and let $Z\subset X$ be of codimension two with a resolution
$0\to E_{1}\to E_{0}\to\cI_{Z}\to0$ by vector bundles of ranks $r$ and $r+1$
and $\ch(\cI_{Z}(b\Theta))=u_{\Theta}+b\,v_{\Theta}$.
\begin{enumerate}[label=\textup{(\roman*)},nosep]
\item $r\ge2$, whatever the bundles. Equivalently, $Z$ is not the zero locus
  of a section of a vector bundle of rank two.
\item If $r=2$ and the Chern classes of $E_{1}$ are polynomials in $\Theta$,
  then $3d-b^{2}-4ab\equiv0\pmod{12}$, where
  $c_{1}(E_{1}(b\Theta))=a\Theta$. At $b=3$ this says $d\equiv3\pmod4$.
\end{enumerate}
\end{proposition}

\begin{proof}
Put $E=E_{1}(b\Theta)$, $G=E_{0}(b\Theta)$, $F=\cI_{Z}(b\Theta)$, so
$c(G)=c(E)\,c(F)$, and Newton's identities give
\[
  c(F)=1+b\Theta+N\Theta^{2}+\tfrac{bN}{3}\Theta^{3}
  +\tfrac{N(b^{2}-3d)}{12}\Theta^{4}.
\]
A vector bundle of rank $s$ has $c_{k}=0$ for $k>s$. At $r=1$, with
$\ell=c_{1}(E)$, the rank two bundle $G$ has
$0=c_{3}(G)=\bigl(\tfrac{bN}{3}\Theta+N\ell\bigr)\Theta^{2}$, so
$\ell=-\tfrac b3\Theta$ by hard Lefschetz, and then
$c_{4}(G)=-\tfrac{N(b^{2}+9d)}{36}\Theta^{4}\ne0$. A section of a rank two
bundle $V$ vanishing in codimension two on $Z$ gives the Koszul resolution
$0\to\det V^{\vee}\to V^{\vee}\to\cI_{Z}\to0$, of rank one. At $r=2$, with
$c_{1}(E)=a\Theta$ and $c_{2}(E)=\tfrac m2\Theta^{2}$, the rank three bundle
$G$ has $0=c_{4}(G)=c_{4}(F)+c_{1}(E)c_{3}(F)+c_{2}(E)c_{2}(F)$, which
integrates to $6m=3d-b^{2}-4ab$, and
$\chi(E)=a^{4}-2a^{2}m+\tfrac12m^{2}\in\ZZ$ makes $m$ even.
\end{proof}

\begin{remark}\label{rem:ranks}
Neither condition restricts $r\ge3$: a bundle of rank four has no vanishing
Chern class on a fourfold, and at $r=3$ the only condition, $c_{4}(E)=0$,
involves $E$ alone \cite[Remark~18.61]{BMB}. From rank three on, Chern
numbers alone decide nothing, and \Cref{cor:examples}(iii) gives data at every
$d\in\{1,3,5,7\}$.
\end{remark}

\begin{status}\label{st:remaining}
Two shapes remain \cite[Remark~18.69]{BMB}: a smooth support at
$d\in\{1,3,5,7\}$ whose Hilbert--Burch resolution has rank $r\ge2$, and
$r\ge3$ at $d=1,5$ in the polynomial case, with one term not projectively
flat and not both terms in the lattice of line bundles; or a singular
Cohen--Macaulay support with a resolution over a bundle that is not
projectively flat. Neither is settled.
\end{status}

\section{The program}\label{sec:program}

\begin{step}[Base point]\label{step:base}
Take any principally polarised abelian fourfold $X$; $d$ is fixed by the
Chern character of the candidate, not by $X$. When $X$ is simple, the
injectivity of $H^{2}(\cO_{X})\to H^{2}(\cO_{S})$ required by
\cite[Lemma~18.62]{BMB} holds for a smooth support automatically
\cite[Lemma~18.43]{BMB}. When $\Hdg^{2}(X)=\QQ\Theta$ and
$\Hdg^{4}(X)=\QQ\Theta^{2}$, as for a very general $X$, all Chern classes are
polynomials in $\Theta$ and \Cref{prop:lowrank}(ii) applies. An action of
$K$ on $X$ is not needed, and it does not help with condition (a)
(\Cref{step:extension}).
\end{step}

\begin{step}[Support]\label{step:support}
Construct a smooth or Cohen--Macaulay surface $Z\subset X$ with
$[Z]=N\Theta^{2}$ and, if smooth, the invariants of \Cref{tab:targets},
with a Hilbert--Burch resolution $0\to E\to G\to\cI_{Z}(3\Theta)\to0$ of
rank $r\ge2$ outside the classes of \Cref{sec:excluded}. The three families
proposed in the first version of this note stand as follows.
\begin{enumerate}[label=(\arabic*),leftmargin=2.2em]
\item Degeneracy loci of maps between Fourier--Mukai transforms. The
  transform of a nondegenerate line bundle on $\widehat X$ is a simple
  semi-homogeneous bundle \cite{Muk78,Muk81}, whose rank on a principally
  polarised fourfold is $q^{4}$ for slope $\tfrac pq\Theta$, never $2$ or $3$;
  a resolution with both terms sums of such bundles is excluded by
  \cite[Theorem~18.66]{BMB}. Transforms of other sheaves are not excluded.
\item Images of $C_{1}\times C_{2}$ or of $C^{(2)}$. A smooth image
  $S\cong C_{1}\times C_{2}$ has $K_{S}^{2}=8\chi(\cO_{S})$, which with
  \Cref{tab:targets} reads $5N=4b^{2}$, impossible at $b=3$ and possible only
  when $5\mid b$. A smooth $S\cong C^{(2)}$ of genus $g$ has
  $\chi(\cO_{S})=\tfrac12(g-1)(g-2)$, never $260$, $288$, $308$ or $320$.
  Singular images must be Cohen--Macaulay and are not excluded.
\item Zero loci of sections of rank two bundles. Excluded by
  \Cref{prop:lowrank}(i), for every bundle, equivariant or not.
\end{enumerate}
\end{step}

\begin{step}[Extension]\label{step:extension}
Verify condition (a) of \Cref{thm:factor} along all ten polarised
directions. If $X$ carries an action of $K$ of signature $(2,2)$, the
$K$-linear polarised directions form a space of dimension $4$, so an
equivariant construction that persists over the Weil family gives (a) along
at most four of the ten directions; the other six destroy the action, and
equivariance says nothing about them. For a smooth support, (a) is
equivalent to the Bloch semiregularity map of $S$ having rank six, the least
value possible \cite[Proposition~18.46]{BMB}.
\end{step}

\begin{step}[The identity \textup{(b)}]\label{step:heisenberg}
For a Cohen--Macaulay support, verify that the images of $A_{\theta}A_{\theta'}$
in $H^{1}(Z,\cExt^{1}(\cI_{Z},\cI_{Z}))$ vanish for all
$\theta,\theta'\in H^{0}(T_{X})$ \cite[Corollary~18.50(ii)]{BMB}. The first
version proposed computing these classes from the Atiyah classes of $E$ and
$G$ restricted to $Z$; that formula is a proposal, it is not in \cite{BMB},
and it has not been carried out for any candidate.
\end{step}

\begin{step}[Propagation]\label{step:propagation}
If \Cref{step:support,step:extension,step:heisenberg} succeed for secant
objects $F_{1},F_{2}$, then $\Phi(F_{1}\boxtimes F_{2})$ on
$X\times\widehat X$ satisfies the weakened criterion
\cite[Theorem~18.29(iv)]{BMB}. With an affirmative answer to
\cite[Question~11.4]{Mar25b}, the semiregularity argument of Buchweitz and
Flenner \cite{BF03} and Pridham \cite{Pri24} would then make
$\mathbf{W}(K,4,\delta_{0})$ algebraic on the eightfolds of the family, and
\cite[Proposition~20.13]{BMB} every $\delta$, for the discriminants reached
and nothing beyond them \cite[Remark~18.69]{BMB}.
\end{step}

\section{The numerical lemma, corrected}\label{sec:lemma}

Write the Chern classes of a class $E$ of rank $r$ as
$c_{k}(E)=m_{k}\Theta^{k}/k!$; since $\Theta^{k}/k!$ is integral and
primitive, integrality means $m_{k}\in\ZZ$, and $m_{4}$ is a number of
points. Put $F=\cI_{Z}(3\Theta)$ with $\ch F=u_{\Theta}+3v_{\Theta}$.

\begin{lemma}\label{lem:whitney}
If $\ch G=\ch E+\ch F$ with $\operatorname{rk}G=r+1$, then $c(G)=c(E)\,c(F)$,
where $c(F)$ has $m_{k}(F)=(3,\,2N,\,6N,\,6N(3-d))$; so $G$ has integral
Chern classes whenever $E$ does. The formulas for $n_{k}=m_{k}(G)$ in the
first version of this note are this product written out. With Bogomolov's
discriminant $\Delta(V)=2\operatorname{rk}V\,c_{2}(V)
-(\operatorname{rk}V-1)c_{1}(V)^{2}$, whose coefficient of $\Theta^{2}$ is
$r\,m_{2}-(r-1)m_{1}^{2}$,
\[
  \Delta(G)-\Delta(E)=\bigl(m_{2}-(m_{1}-3)^{2}+18+(r+1)d\bigr)\Theta^{2}.
\]
\end{lemma}

The first version displayed $\operatorname{rk}V\,c_{2}(V)
-(\operatorname{rk}V-1)c_{1}(V)^{2}$, which differs from Bogomolov's in the
factor $2$ on $c_{2}$, while its values, such as $\Delta(E)=4\Theta^{2}$ for
$c(E)=1+\Theta^{2}$ in rank two, are Bogomolov's. What the lemma does not
see is the vanishing of $c_{k}(E)$ for $k>r$ and of $c_{k}(G)$ for $k>r+1$,
and this is the whole constraint at $r\le2$.

\begin{corollary}\label{cor:examples}
Let $b=3$ and $d\in\{1,3,5,7\}$, and suppose the Chern classes of $E$ are
polynomials in $\Theta$.
\begin{enumerate}[label=\textup{(\roman*)},nosep]
\item Rank two data exist exactly at $d=3$ and $d=7$. With
  $c_{1}(E)=a\Theta$ they have $\Delta(E)=\bigl(d+1-(a+2)^{2}\bigr)\Theta^{2}$
  and $\Delta(G)=\tfrac12\bigl(9(d+1)-4a^{2}\bigr)\Theta^{2}$.
\item The example of the first version, $c(E)=1+\Theta^{2}$ in rank two, has
  $c_{4}(G)=180,144,84,0$ points at $d=1,3,5,7$, so it is a resolution datum
  only at $d=7$. With the conditions of the first version,
  $\Delta(E)>0$, $\Delta(G)>0$, $\mu(E)<\mu(G)$ and
  $\chi(E^{\vee}\otimes G)\ge1$, rank two survives only at $d=7$, for
  $a\in\{-2,-1,0\}$; at $d=3$ the two discriminants are positive only for
  $a\in\{-2,-1\}$, where $\chi(E^{\vee}\otimes G)=-78$ and $-7$.
\item In rank three, $c(E)=1+\Theta^{2}$ meets all four conditions at every
  $d\in\{1,3,5,7\}$, with $\chi(E^{\vee}\otimes G)=53,101,173,269$.
\end{enumerate}
\end{corollary}

\begin{proof}
At $r=2$ the relation of \Cref{prop:lowrank} gives
$m_{2}=\tfrac12(d-3)-2a$, from which the two discriminants follow. The
remaining values are finite computations, checked exactly in
\Cref{sec:checks}; the closed form of $\chi(E^{\vee}\otimes G)$ at $r=2$ is
$\tfrac18\bigl(24a^{4}+64a^{3}-88a^{2}d+104a^{2}-64ad+192a+57d^{2}-174d
+153\bigr)$.
\end{proof}

The four conditions of the first version are heuristics: Bogomolov's
inequality bounds $\Delta$ below for a stable bundle, and
$\chi(E^{\vee}\otimes G)\ge1$ neither implies nor is implied by the existence
of a map $E\to G$. Chern numbers therefore decide ranks one and two and say
nothing from rank three on.

\section{The toy computation}\label{sec:toy}

The first version searched unions of coordinate abelian surfaces
$\{x_{i}=c_{i},x_{j}=c_{j}\}$ in $X=E_{1}\times E_{2}\times E_{3}\times E_{4}$
with the product polarisation, at $b=1$, for configurations with the Chern
character of a secant support, Cohen--Macaulay everywhere, and with a
combinatorial condition proposed there as necessary for (a). The programs are
in \texttt{toy/}. A rerun of the search at $b=1$, $d=1$ with forty restarts
of $60000$ steps found an exact configuration on four fibres per factor,
twelve surfaces with the required pair and triple counts and
$\chi(\cO_{Z})=4$, and none on three fibres per factor, where the best
configuration missed by one in a triple count. At $d=3$ twenty restarts on
each size found none; the best configurations missed the targets by a total
of $31$ and $16$. The
symbolic powers of the union of the six coordinate surfaces through a point
give no admissible Chern character for $1\le m\le20$.

This says nothing about \Cref{step:support} on a simple fourfold: $b=1$ is
not Markman's normalisation, $X$ is not simple and $\NS(X)$ has rank at
least four, so neither \cite[Lemma~18.43]{BMB} nor \Cref{prop:lowrank}(ii)
applies, and abelian surfaces do not survive a deformation of $X$ to a
simple fourfold.

\section{The checks}\label{sec:checks}

The program \texttt{route\_checks.py} performs nineteen checks in exact
arithmetic: the invariants of \Cref{tab:targets} and the class
$[S]=N\Theta^{2}$; \Cref{lem:whitney} and the formulas of the first version
on random data; the Bogomolov normalisation; the example and the three parts
of \Cref{cor:examples}; the value $c_{4}(G)=-6N(1+d)$ points at $r=1$,
$b=3$; the conditions on products and symmetric squares of curves; the rank
$q^{4}$ of a simple semi-homogeneous bundle; and the dimension $pq$ of the
$K$-linear polarised directions for signature $(p,q)$. Item (LXX) of
\cite[Appendix~D]{BMB} checks \Cref{prop:lowrank}.

\section{Changes from the first version}\label{sec:changes}

\begin{enumerate}[label=(\arabic*),leftmargin=2.2em]
\item The target. The Hodge conjecture for abelian fourfolds of Weil type is
  Markman's theorem; the route aims at $\mathbf{W}(K,4,\delta)$ on eightfolds
  and needs \cite[Question~11.4]{Mar25b}.
\item The closure framing. (F1) is needed only as $\mathrm{P2}_{\mathrm{split}}$,
  for CM fields of degree at least four; the secant route lies in no minimal
  sufficient set; \cite[Proposition~20.35]{BMB} concerns (F3), not (F3$'$).
\item The base point needs no $K$-action, and a $K$-action gives condition (a)
  along at most four of ten directions.
\item The class of a smooth support is $N\Theta^{2}$, not $N\Theta^{2}/2$.
\item The local statement is one implication, and the global condition is a
  family of classes indexed by pairs $\theta,\theta'$.
\item Family (3) is excluded, family (1) is excluded for transforms of line
  bundles, and smooth products and symmetric squares of curves are excluded
  at $b=3$.
\item The numerical lemma omits the vanishing above the rank; its example
  works only at $d=7$; its displayed discriminant lacks a factor $2$.
\item The bibliography is that of \cite{BMB}: \cite{Mar25b} is ``Secant
  sheaves and Weil classes on abelian varieties'', and Pridham's paper is
  \cite{Pri24}.
\end{enumerate}

\begin{thebibliography}{Mar25b}

\bibitem[BMB]{BMB}
D. Bhattacharjee, P. Mandal and U. Bhattacharya,
\emph{Density of the Algebraic Locus of Weil Classes on Abelian Varieties},
manuscript, 2026; code and sources:
\href{https://doi.org/10.5281/zenodo.23038895}{\texttt{doi:10.5281/zenodo.23038895}}.
Theorem numbers refer to the version of 29 September 2026.

\bibitem[BF03]{BF03}
R.-O. Buchweitz and H. Flenner,
\emph{A semiregularity map for modules and applications to deformations},
Compositio Math. \textbf{137} (2003), no.~2, 135--210.
\href{https://doi.org/10.1023/A:1023999012081}{\texttt{doi:10.1023/A:1023999012081}}.

\bibitem[Mar25a]{Mar25a}
E. Markman,
\emph{Cycles on abelian $2n$-folds of Weil type from secant sheaves on
abelian $n$-folds}, preprint, arXiv:2502.03415 (2025).
\href{https://doi.org/10.48550/arXiv.2502.03415}{\texttt{doi:10.48550/arXiv.2502.03415}}.

\bibitem[Mar25b]{Mar25b}
E. Markman,
\emph{Secant sheaves and Weil classes on abelian varieties},
preprint, arXiv:2509.23403 (2025).
\href{https://doi.org/10.48550/arXiv.2509.23403}{\texttt{doi:10.48550/arXiv.2509.23403}}.

\bibitem[Muk78]{Muk78}
S. Mukai,
\emph{Semi-homogeneous vector bundles on an abelian variety},
J. Math. Kyoto Univ. \textbf{18} (1978), no.~2, 239--272.
\href{https://doi.org/10.1215/kjm/1250522574}{\texttt{doi:10.1215/kjm/1250522574}}.

\bibitem[Muk81]{Muk81}
S. Mukai,
\emph{Duality between $D(X)$ and $D(\hat{X})$ with its application to Picard
sheaves},
Nagoya Math. J. \textbf{81} (1981), 153--175.
\href{https://doi.org/10.1017/S002776300001922X}{\texttt{doi:10.1017/S002776300001922X}}.

\bibitem[Pri24]{Pri24}
J. P. Pridham,
\emph{Semiregularity as a consequence of Goodwillie's theorem},
Forum Math. Sigma \textbf{12} (2024), e126. arXiv:1208.3111.
\href{https://doi.org/10.1017/fms.2024.132}{\texttt{doi:10.1017/fms.2024.132}}.

\bibitem[Sch98]{Sch98}
C. Schoen,
\emph{Addendum to: Hodge classes on self-products of a variety with an
automorphism},
Compositio Math. \textbf{114} (1998), no.~3, 329--336.
\href{https://doi.org/10.1023/A:1000566205021}{\texttt{doi:10.1023/A:1000566205021}}.

\end{thebibliography}

\end{document}
